\section{Environment}
\label{section:environment}

The environment is implemented with \texttt{poke-env}, connected to a local Pokemon Showdown server.
In our code, \texttt{PokeRLEnv} extends \texttt{SinglesEnv} (multi-agent), and is wrapped by
\texttt{MaskableSingleAgentWrapper} to expose a standard single-agent \texttt{gymnasium.Env}
interface for Stable-Baselines3.

At each decision step, \texttt{poke-env} provides a \texttt{Battle} object (state snapshot of the
current turn). Our observation pipeline is:
\[
\texttt{Battle state} \;\longrightarrow\; \texttt{embed\_battle(battle)} \;\longrightarrow\; \mathbf{s}_t \in \mathbb{R}^{490}.
\]

The \texttt{Battle} object is not a flat vector: it is a structured Python object with many fields.
Our encoder reads only the fields needed for policy learning:
\begin{itemize}
	\item active Pokemon on both sides (HP fraction, types, boosts, status),
	\item known moves of the active Pokemon,
	\item revealed team members on both sides (bench information),
	\item side conditions (entry hazards and screens),
	\item global effects (weather and field effects),
	\item availability flags for battle gimmicks (Dynamax, Mega, Z-Move, Tera).
\end{itemize}

This is still a partially observable setting: opponent information is only available when revealed
during battle (e.g., observed Pokemon and seen moves).

The state vector is built by concatenating feature blocks in a fixed order:
\[
\mathbf{s}_t = [\mathbf{s}^{ally\_active},\; \mathbf{s}^{ally\_moves},\; \mathbf{s}^{opp\_active},\; \mathbf{s}^{opp\_moves},\; \mathbf{s}^{ally\_bench},\; \mathbf{s}^{opp\_bench},\; \mathbf{s}^{side},\; \mathbf{s}^{weather},\; \mathbf{s}^{field},\; \mathbf{s}^{gimmick}].
\]

\begin{itemize}
	\item \textbf{Ally active Pokemon (33)}: HP fraction (1), type multi-hot over 18 types,
	7 normalized stat boosts (\(-6\) to \(+6\) scaled to \([-1,1]\)), and status one-hot (7).

	\item \textbf{Ally active moves (92)}: 4 fixed slots, each encoded with 23 values:
	move type one-hot (18), base power/250, accuracy/100, category bits (physical/special),
	and normalized PP \(\texttt{current\_pp}/\texttt{max\_pp}\).

	\item \textbf{Opponent active Pokemon (33)}: same encoding as ally active Pokemon.

	\item \textbf{Known opponent moves (220)}: up to 10 known moves, 22 values each
	(same as above \emph{without} PP, because opponent PP is not exposed by the server).
	Unknown slots are zero-padded.

	\item \textbf{Ally bench (40)} and \textbf{opponent bench (40)}:
	5 slots per side, each slot has HP fraction (1) + status one-hot (7), for 8 values per slot.
	Missing/unrevealed slots are zero-padded.

	\item \textbf{Side conditions (14)}:
	for each side, 4 hazard values (Stealth Rock, Spikes layers, Toxic Spikes layers, Sticky Web)
	and 3 binary screen flags (Reflect, Light Screen, Aurora Veil).

	\item \textbf{Weather (8)}: one-hot over the supported weather states.

	\item \textbf{Field effects (6)}: one-hot over terrain/global field effects
	(Electric/Grassy/Misty/Psychic Terrain, Trick Room, Gravity).

	\item \textbf{Gimmick availability (4)}:
	\([\texttt{can\_dynamax},\texttt{can\_mega\_evolve},\texttt{can\_tera},\texttt{can\_z\_move}]\).
\end{itemize}

Therefore, the final observation size is:
\[
33 + 92 + 33 + 220 + 40 + 40 + 14 + 8 + 6 + 4 = 490.
\]

All features are cast to \texttt{float32}. Continuous values are normalized (HP, boosts,
power, accuracy, PP), categorical values are encoded as one-hot/multi-hot vectors, and variable-length
information (moves, bench, revealed opponent data) is converted to fixed-size blocks by zero-padding.
